In Cosmology, observations of the spatial distribution of galaxies are crucial
to understanding the evolution and expansion of the Universe. Numerous
telescopes conduct surveys to produce catalogues of galaxies, which are then
subjected to statistical analysis. Usually, galaxies form in groups, which
together are part of a cluster, surrounded by a very large, spherical
dark-matter halo. Most of the matter is concentrated in the centre, thus there
is a high probability of having a large-mass galaxy at the centre of the halo.

In this problem, we will assume that all halos have one central galaxy and
several groups of satellite galaxies, and use statistical analysis to draw
inferences about the spatial distribution of galaxies. In Figure 1, you can
see six different possible probability densities (A--F) for distances between
galaxies. On the X axis is the separation distance between galaxies ($a$), and
on the Y axis is the probability density ($f_p(a)$). The probability of the
distance between a pair of galaxies being in the range $a$ to $a + \Delta a$
is given by
\[ P(a) = f_p(a)\,\Delta a \]
Note that $f_p(a)$ may take values greater than 1 in some cases.

\ioaafig{2022_DA_2-f5.pdf}
% densities f_p(a) vs separation a (source paper, Data Analysis page 3 of 7,
% "Figure 1: Probability density of all possible galaxy pairs")

\medskip
\textbf{Part 1 (60\,pt)}

Table 1 gives the observational data of a single dark matter halo in the form
of Cartesian coordinates of 16 satellite groups and the coordinates of the
central galaxy, which is given in the first row. The $i^{\mathrm{th}}$ group
consists of $N_i$ galaxies. The origin of this coordinate system is the Earth.
For simplicity we assume that a group of galaxies means a collection of
galaxies with almost the same coordinates.

\begin{center}
\begin{tabular}{ccccc}
$i$ & $x$ (Mpc) & $y$ (Mpc) & $z$ (Mpc) & $N_i$ \\
\hline
C  & 8.300 & 6.200 &  1.100 & Central \\
\hline
1  & 8.401 & 6.309 &  1.394 & 291 \\
2  & 9.883 & 7.189 &  1.506 & 170 \\
3  & 8.883 & 6.413 &  2.226 & 253 \\
4  & 8.444 & 6.569 &  2.439 & 8   \\
5  & 7.782 & 6.048 & $-0.358$ & 72  \\
6  & 8.780 & 6.305 &  2.463 & 120 \\
7  & 7.990 & 5.881 &  0.532 & 302 \\
8  & 8.540 & 6.388 &  0.369 & 52  \\
9  & 7.975 & 6.030 &  1.216 & 28  \\
10 & 7.072 & 5.431 &  0.488 & 40  \\
11 & 8.037 & 5.965 &  0.615 & 72  \\
12 & 8.681 & 6.483 &  0.734 & 379 \\
13 & 7.115 & 5.259 &  1.403 & 82  \\
14 & 9.587 & 7.130 &  0.322 & 62  \\
15 & 8.193 & 6.104 &  0.915 & 305 \\
16 & 7.613 & 5.847 &  1.666 & 293 \\
\end{tabular}
\end{center}

\begin{description}
\item[(a)] The distribution of the satellite groups in the $X-Z$ plane and
  $Y-Z$ plane has been plotted for you. Plot the distribution in the $X-Y$
  plane on millimeter paper. \hfill(6\,pt)

  \ioaafig{2022_DA_2-f2.pdf}

\item[(b)] Given that this cluster of galaxies is shaped like a disk and our
  line of sight is in the plane of the disk, find the unit vector in the
  direction of the normal to the disk. \hfill(4\,pt)
\item[(c)] Let us define a new coordinate system centered on the central
  galaxy, defining our line of sight from the Earth to the central galaxy as
  the $X'$ axis and the normal to the disk as the $Z'$ axis. Find the unit
  vectors for this coordinate system. \hfill(2\,pt)
\item[(d)] Convert the given Cartesian coordinates to the new $X'Y'Z'$
  coordinates. Calculate $X'$, $Y'$, $Z'$. \hfill(24\,pt)
\item[(e)] Write the equations necessary to convert the new $X'Y'Z'$
  Cartesian coordinates into spherical coordinates centered on the central
  galaxy. The $\theta$ coordinate has been already calculated for you below.
  Calculate $R$ and $\phi$. \hfill(11\,pt)

  \begin{center}
  \begin{tabular}{cc}
  $i$ & $\theta$ \\
  \hline
  1  & 1.652 \\
  2  & 1.490 \\
  3  & 1.432 \\
  4  & 1.721 \\
  5  & 1.692 \\
  6  & 1.430 \\
  7  & 1.474 \\
  8  & 1.580 \\
  9  & 1.723 \\
  10 & 1.646 \\
  11 & 1.519 \\
  12 & 1.569 \\
  13 & 1.542 \\
  14 & 1.557 \\
  15 & 1.516 \\
  16 & 1.705 \\
  \end{tabular}
  \end{center}

  \emph{Note:} in spherical coordinates, one uses the coordinates $R$,
  $\theta$ and $\phi$, where $R$ is the radial distance, $\phi$ is the
  azimuthal angle measured anti-clockwise from the $X'$ axis (from 0 to
  $2\pi$) and $\theta$ is the polar angle (from 0 to $\pi$), i.e.\ the angle
  relative to the $Z'$ axis.
\item[(f)] Plot a histogram (on millimeter paper) showing the number of
  satellite galaxies at different radial distances from the central galaxy.
  Use 8 equally sized bins for radial distances, with 0.2\,Mpc to 2\,Mpc as
  the total range of the histogram. \hfill(9\,pt)
\item[(g)] The histogram must also be converted into a probability density. On
  the right hand side of the vertical axis of the histogram you plotted, add
  markings for the re-scaled vertical axis appropriate for the probability
  density. You will use this probability density data in part 3. \hfill(4\,pt)
\end{description}

\medskip
\textbf{Part 2 (17\,pt)}

The coordinates in the Table above are determined using distances obtained by
the redshift method. It is well-known that while measuring distance to any
galaxy with the redshift method, errors might occur because of the orbital
motion of the galaxies. Now, we will try to correct for these errors. Assume
that the central galaxy has an effective mass of $1 \times 10^{10}\,M_\odot$,
all satellite galaxies move in circular orbits around the central galaxy and
interactions between any two satellite galaxies can be ignored.

\begin{description}
\item[(a)] Use the data in Table 1 of the question sheet to obtain a histogram
  of number of galaxies against the $\phi$ azimuthal angle (take the width of
  the bins as $60^\circ$). The angle distribution is expected to be uniform,
  but it is not. Calculate the standard deviation for this data. \hfill(5\,pt)
\item[(b)] The two bins centered along the $Y'$ axis show an unusually low
  number of galaxies as the orbital velocities of satellite galaxies in these
  bins are almost parallel to our line of sight and hence there is maximum
  error in their distance estimated by the redshift method. Derive an
  expression for the correction in each azimuthal angle. \hfill(4\,pt)

  \emph{Note:} You may assume that the distance from the central galaxy ($R$)
  doesn't change significantly because of this correction.
\item[(c)] Calculate the corrected value of the azimuthal angles for each data
  row and recalculate the histogram and its standard deviation of its values.
  \hfill(8\,pt)
\end{description}

\medskip
\textbf{Part 3 (28\,pt)}

So far, we have only calculated the central galaxy -- satellite galaxy
histogram for a single halo. It is time to move on to calculating other
relations. The same observations and data analysis has been carried out for
different halos and their central galaxies. Suppose we have 500 halos, thus,
500 central galaxies ($N_C = 500$). Take that the average number of satellite
galaxies in each halo is 1000 ($N_S = 1000$).

\begin{description}
\item[a)] Calculate the following quantities:
  \begin{description}
  \item[(1)] $N_{CC}$, total number of pairs of central galaxies.
  \item[(2)] $N_{CS}^*$, total number of central galaxy -- satellite pairs
    such that the galaxies in the pair belong to the same halo.
  \item[(3)] $N_{CS}$, total number of central galaxy -- satellite pairs such
    that the galaxies in the pair belong to two separate halos.
  \item[(4)] $N_{SS}^*$, total number of satellite -- satellite galaxy pairs
    such that the galaxies in the pair belong to the same halo.
  \item[(5)] $N_{SS}$, total number of satellite -- satellite galaxy pairs
    such that the galaxies in the pair belong to two separate halos.
  \end{description}

  The table below gives the probability density values for different distances
  between central galaxies in different dark matter halos (CC). You should
  note this data is not real and doesn't represent the results of the real
  observations. Dark matter halos do not intersect. But for simplicity we will
  use this data.

  \begin{center}
  \begin{tabular}{cc}
  $R$ (Mpc) & $\rho_{CC}$ \\
  \hline
  0.3125 & 0.406 \\
  0.5375 & 0.808 \\
  0.7625 & 1.341 \\
  0.9875 & 1.134 \\
  1.2125 & 0.479 \\
  1.4375 & 0.143 \\
  1.6625 & 0.073 \\
  1.8875 & 0.060 \\
  \end{tabular}
  \end{center}

  We can calculate all probability densities using the given Center-Center
  (CC) data and the same halo Center-Satellite ($CS^*$) data that we have
  calculated in part 1. We use the following relations:
  \[ \rho_{SS}^*(R) = c_1 R\left(\rho_{CS}^*(R)\right)^2 \]
  \[ \rho_{CS}(R) = c_2 R\left(\rho_{CS}^*(R)\right)^2
     \sqrt{\rho_{CC}(R) + 1} \]
  \[ \rho_{SS}(R) = c_3 R\left(\rho_{CS}^*(R)\left(\rho_{CC}(R)
     - 5\right)\right)^2 \]
  where $c_1$, $c_2$ and $c_3$ are normalization constants and should be found
  for each probability density.
\item[b)] Calculate normalized $\rho_{SS}^*(R)$, $\rho_{CS}(R)$,
  $\rho_{SS}(R)$. \hfill(15\,pt)
\item[c)] Calculate the final probability density of all galaxy pairs, that
  combines all the distributions above with appropriate weights. Which one of
  the 6 graphs on Fig.~1 best represents this probability density?
  \hfill(8\,pt)
\end{description}
